% ~1.25 pages. Ported 2026-11-21 from ../sec_system_model_waveform.tex (cgan/, 30 dB) and updated to final/
% (README Sec. 3.4 "Final experiment", 2.7): one jammer (no team: every learned jammer in final/ is N_J = 1);
% the victim link fades and the noise level is uncertain (the parked sentences of ../sec_setup_operating_point.tex);
% energy after the matched filter, plus a gain-aware variant; the noise LRT is dropped (final/ has none).
% Symbols: g = RRC pulse; rho = fading amplitude (NOT g, which the old parked text used); nu = noise factor;
% T = a detector's statistic; tau_alpha its threshold. Style (ADM): one sentence per line, ~\cite before punctuation.
\section{System and Threat Model}\label{sec:system}
We consider a single-carrier link on which a legitimate node transmits QPSK to a second legitimate node, which also monitors the link for interference.
A jammer seeks to raise the receiver's error rate while drawing as few detections as possible.
The transmitter, the receiver and its detectors form one fixed target; only the jammer adapts.

\subsection{Legitimate Link}\label{sec:link}
The transmitter sends frames of $N$ unit-energy QPSK symbols $s[k]\in\{\tfrac{1}{\sqrt2}(\pm1\pm\mathrm{j})\}$, drawn independently and uniformly.
Each frame is upsampled by $L$ samples per symbol and shaped by a unit-energy root-raised-cosine pulse $g$, giving the transmit waveform
\begin{equation}
x[n]=\sum_{k=1}^{N} s[k]\,g[n-kL].
\label{eq:tx}
\end{equation}
The receiver applies the matched filter, samples it at the symbol instants to obtain $z_k$, and makes a hard decision on each axis.
Because $g$ is real and even, the matched filter is $g$ itself, and because $g$ is root-Nyquist the cascade $g*g$ is free of inter-symbol interference up to truncation.

\subsection{Channel and Received Signal}\label{sec:channel}
The received frame is
\begin{equation}
r[n]=\rho\,x[n]+c[n]+w[n],\qquad w[n]\sim\mathcal{CN}(0,\nu N_0),
\label{eq:rx}
\end{equation}
where $c$ is the jammer's waveform as it arrives at the receiver.
The amplitude $\rho=|h|$, with $h$ Rician of factor $K$ and $\mathbb{E}[\rho^2]=1$, models small-scale fading of the victim's own link, and the factor $\nu>0$, log-normal with $\mathbb{E}[\nu]=1$ and standard deviation $\sigma_N$ in decibels, models a noise level the receiver does not know exactly.
Both are drawn once per frame, and neither is known to the detectors.
The receiver is synchronized to its transmitter and knows the link's phase, so no equalization is needed; $K\to\infty$ and $\sigma_N=0$ recover the ideal link.
% final/link.py fading_gain (h = sqrt(K/(K+1)) + sqrt(1/(K+1)) CN(0,1)) and noise_scale (unit-mean log-normal).
All powers are means per sample at the receiver input: with $P_{\mathrm{rx}}$ the mean received signal power, the signal-to-noise ratio is $P_{\mathrm{rx}}/N_0$ per sample, so that $E_s/N_0=L\,P_{\mathrm{rx}}/N_0$.
The jammer is characterized by its jamming-to-signal ratio (JSR), its received power relative to $P_{\mathrm{rx}}$.
Every realizable jammer is asynchronous: its timing offset is uniform over one symbol and its carrier phase uniform over $[0,2\pi)$, both drawn per frame.
% Only the genie (T0) and the random push (Sec. method) arrive synchronously.

\subsection{Attacker Model}\label{sec:attacker}
Every attacker, closed-form or learned, is a conditional law $c\sim\pi(\cdot\mid\xi)$, where $\xi$ is the information available when it forms its signal; attackers differ only in $\xi$ and in the family $\pi$ ranges over.
The power budget is a hard constraint: each frame is scaled so that the jammer's received power over the victim's $NL$-sample frame meets its target,
\begin{equation}
\frac{1}{NL}\sum_{n}|c[n]|^2=\mathrm{JSR}\cdot P_{\mathrm{rx}},
\label{eq:power}
\end{equation}
so no attacker can trade effectiveness against a soft power penalty.
Three information tiers define $\xi$: a genie (T0) knows the victim's current symbols and arrives synchronously; an informed attacker (T1) is asynchronous and knows the victim's modulation, pulse and symbol rate but not its symbols; a blind attacker (T2) is asynchronous and knows none of these.

The learned jammer is generative: a neural network $G$ maps a latent vector $\mathbf{z}\sim\mathcal{N}(\mathbf{0},\mathbf{I})$ to a segment $G(\mathbf{z})\in\mathbb{C}^{M}$ of $M$ complex baseband samples.
A jammer frame is a stream of independent segments, delayed and rotated asynchronously and scaled to its budget,
\begin{equation}
u=\bigl[G(\mathbf{z}_1),G(\mathbf{z}_2),\ldots\bigr],\qquad c[n]=a\,e^{\mathrm{j}\phi}\,u[n-\tau],
\label{eq:gan}
\end{equation}
with $(\tau,\phi)$ the offset and phase of Sec.~\ref{sec:channel} and $a>0$ the amplitude that meets~\eqref{eq:power}.
Its weights are learned in simulation (Sec.~\ref{sec:method}), which places it in tier T1; it never sees the victim's symbols.
The classical and reference jammers are defined in Sec.~\ref{sec:method}.

\subsection{Detector Model}\label{sec:detector}
Detection is at the receiver, on one frame at a time, without channel state information or labels at run time.
Each detector reduces the frame to a scalar statistic $T$ built so that larger is more suspicious, and decides alone.
Its threshold $\tau_\alpha$ is the $(1-\alpha)$-quantile of $T$ over clean frames of the environment it is deployed in, including the fading and noise-level draws it cannot know, so the false-alarm rate stays at $\alpha$ and an impairment raises the threshold instead.
The detection probability at a given JSR is
\begin{equation}
P_{\mathrm{det}}=\Pr\bigl[T(r)>\tau_\alpha\mid\text{jammer at JSR}\bigr],\qquad \Pr\bigl[T(r)>\tau_\alpha\mid\text{no jammer}\bigr]=\alpha,
\label{eq:pdet}
\end{equation}
where both probabilities run over everything random in the environment: bits, noise, channel, and the jammer's latent input, delay and phase.
It is a per-frame probability, not the probability that the jammer is ever caught.

We deploy three detectors.
The energy detector~\cite{urkowitz1967energy} thresholds the mean received symbol energy after the matched filter, $T_{\mathrm{E}}=\frac1N\sum_k|z_k|^2$, which excludes the noise outside the signal's band.
Its one-sided form flags excess energy; a two-sided form, which also flags a deficit, is evaluated but not trained against.
Under fading we also evaluate a gain-aware variant that is handed $\rho$ and removes its effect on the signal term; it bounds what a receiver that estimates its gain could achieve.
% energy_csi subtracts (rho^2 - 1) c0^2 (final/detectors.py); not a deployable detector, a bracket.
The kurtosis detector thresholds the deviation of $\kappa(r)=\mathbb{E}|r|^4/(\mathbb{E}|r|^2)^2$ from its clean mean, in both directions; $\kappa$ is scale-free, so it sees the shape of the amplitude distribution, which energy cannot.
The learned detector is the spectrogram classifier of Li et al.~\cite{li2022jamming_ofdm}: a convolutional network applied to the two-sided complex spectrogram of $r$ on a fixed decibel scale, retrained on this link with their jammer types.
% fixed dB scale, not a per-image contrast stretch: the stretch erases a jammer that lifts the band evenly (README A.4).
Its statistic is the logit difference of its two classes.

\subsection{Threat Model}\label{sec:threat}
The detectors are calibrated, and the classifier trained, before any attack; they are then frozen, and a trained jammer is frozen at deployment, so any adaptation is offline.
A white-box jammer differentiates the deployed detector's statistic, including the classifier's weights; this is the attacker's upper bound, and every generator in Sec.~\ref{sec:res-matched} is trained with it.
A grey-box jammer knows the classifier's architecture and training recipe but not its weights, and trains against a surrogate of its own instead~\cite{papernot2017practical}.
For the energy and kurtosis detectors the two coincide.
A learned jammer sees the victim's error rate only in simulation, during training; a deployed jammer cannot measure it.

\subsection{Problem Formulation}\label{sec:objective}
The jammer maximizes the error rate it induces, penalized by detection,
\begin{equation}
\max_{\pi}\;\mathbb{E}[\mathrm{BER}]-\beta\,P_{\mathrm{det}}\quad\text{s.t.~\eqref{eq:power}},
\label{eq:problem}
\end{equation}
where $\beta\ge0$ weighs detection against damage and power enters only through the hard constraint.
Since each $\beta$ reaches only the upper concave envelope of the achievable set, the object of study is the whole trade-off between damage and $P_{\mathrm{det}}$ for each detector, compared at matched damage rather than at matched power or configuration.
